\section{Discussion and Threats to Validity}

The central finding is a protocol effect, not an implementation bug: two
independent codebases, one built to match Farid et al.'s own tools (Weka
J48, a textbook NB on nominal columns) and one built independently
(sklearn entropy trees, one-hot attributes, a mixed-likelihood NB), agree
that fitting Algorithm~1's filter on the whole dataset before
cross-validation inflates its reported accuracy by double digits, while
fitting Algorithm~2's attribute weights the same way inflates almost
nothing. That agreement across two different attribute encodings and two
different tree implementations is evidence the effect belongs to the
algorithm's structure (Algorithm~1 removes rows that then vanish from
every test fold; Algorithm~2 only removes columns, which are never
test-fold-specific) rather than to a coding choice in either
implementation.

Several limitations bound what this replication can claim. The R1 tree is
Weka's J48, a close match to Quinlan's C4.5~\cite{quinlan1993c45}, but E3a
and E9 use scikit-learn's~\cite{pedregosa2011scikit} CART-style entropy
tree on one-hot attributes, which is not C4.5 and does not test a nominal
attribute the way the original algorithms describe; the E3a/E9 numbers
should be read as a leakage audit on a related but distinct
implementation, not as a second attempt to match the paper's Table 7--11
numbers. Ten datasets is a small sample for any between-dataset
statistical test; every Wilcoxon test in this paper has $n=10$ pairs at
most, which is why several comparisons with visually large mean
differences (for instance Algorithm~1 vs.\ C4.5 under protocol B, median
$-3.71$, raw $p=0.055$) do not survive Holm correction, and why we report
these as directional rather than as confirmed effects. All datasets come
from the paper's original 2014 benchmark, several of them (contact lenses,
iris, vote) small and decades old; nothing here speaks to how either
algorithm behaves on larger or more recent data outside this
benchmark. The seed counts (3 for R1 and E3a, 5 for E9) are smaller than
the 10 originally planned for this line of experiments, and the results
should be read as preliminary; increasing seed count would narrow the
per-arm standard deviations shown in Table~\ref{tab:e9} without changing
which comparisons are paired by dataset for the Wilcoxon tests. No
per-fold accuracy is retained in the result files produced by this
project's code, only fold-averaged, seed-averaged accuracy per dataset and
arm; a corrected resampled $t$-test that accounts for the dependence
between folds from the same dataset was therefore not computed, and we
report only the between-dataset Wilcoxon test, following
Dem\v{s}ar~\cite{demsar2006statistical}, rather than claim a test this
data cannot support.

The NSL-KDD C4.5 discrepancy (99.47\% here vs.\ 71.11\% reported) remains
unresolved and is reported plainly rather than explained away. NSL-KDD is
also the dataset with the largest number of attributes and instances in
the benchmark, and a leakage effect specific to intrusion-detection
preprocessing has been documented on it
elsewhere~\cite{bouke2023empirical}; whether the paper's 71.11\% reflects
a different data split, a different feature set, or a different intrusion
detection preprocessing step is not something this replication can
determine from the published paper alone.

This audit tests one filter design (a single NB classifier judging its
own training data) and one weighting design (inverse-square-root tree
depth); it does not test whether a consensus or ensemble
filter~\cite{brodley1999identifying}, a confidence-thresholded
filter~\cite{northcutt2021confident}, or a correction-based
alternative~\cite{jiang2024correction} would recover some of the honest
accuracy that Algorithm~1's hard deletion gives up, nor whether a
multi-view attribute weighting scheme~\cite{zhang2023multiview} would
close Algorithm~2's small remaining gap to the paper. Those are separate,
un-run experiments; this paper's claim is limited to what Algorithms~1 and
2, as specified in the 2014 paper, do under leakage-safe evaluation on
these ten datasets.
